\section{从摩尔定律时代到后登纳德时代}

\subsection{两条经常被混淆的缩放规律}

现代并行计算的硬件动因必须从两条历史趋势区分开来。

\begin{table}[H]
\centering
\caption{摩尔定律与登纳德缩放的作用区分}
\begin{tabularx}{\textwidth}{L{0.24\textwidth}Y Y}
\toprule
\textbf{概念} & \textbf{核心含义} & \textbf{对性能的历史影响} \\
\midrule
\term{摩尔定律}{Moore's Law} & 制造特征尺寸不断缩小，使单位面积可容纳的晶体管数量约每 18--24 个月翻倍。它是经验性“定律”，也可视为产业目标。 & 提供更高的元件密度，使芯片能够容纳更多功能单元、缓存和控制逻辑。它本身并不等价于单线程性能翻倍。 \\
\term{登纳德缩放}{Dennard Scaling} & 随器件尺寸缩小，晶体管能够在更高时钟频率下工作，是推动频率上升的物理基础。 & 使单个处理器核心能够更快执行同一指令流，是早期“相同代码自动变快”的重要来源。 \\
\bottomrule
\end{tabularx}
\end{table}

大众叙述常把处理器性能增长统称为“摩尔定律”，但历史上的快速性能增长实际来自摩尔定律与登纳德缩放的共同作用。晶体管数量增加提供了更复杂的处理器结构，时钟频率提高则直接缩短许多指令的执行时间。两者叠加后，通用计算机在约三十年内表现出大致每 18 个月翻倍的性能趋势。

\subsection{“自动性能增长”对软件的影响}

在这一时期，软件常能在不改变算法与并行结构的情况下获得显著加速。一个今天运行过慢的串行程序，经过一到两个处理器代际后，只需重新编译或迁移到新机器，就可能明显变快。由此形成了三项长期软件预期。

\subsubsection{同一份软件随硬件代际自动加速}

程序不需要显式使用更多线程，也能受益于更高频率、更好的流水线与更强的单核结构。在这种历史环境中，“代码不够快，等待两年”一度可以成立；它不是现代系统中的可靠开发策略。

某些程序较早暴露了这一策略的边界。若程序主要受\term{内存访问延迟}{memory-access latency}限制，例如部分稀疏、非结构化或图算法，处理器算术速度提高并不会按相同比例缩短运行时间。此类程序需要的数据访问不规则，缓存难以持续命中，因此内存系统可能先成为瓶颈。

\subsubsection{内存从稀缺资源变为多数应用中的充足资源}

芯片密度增加也扩大了可用内存容量。早期程序常需要显式把数据分块，在磁盘与主存之间来回搬运；随着系统能够提供数百吉字节乃至更大容量，许多应用可以把完整数据集常驻内存。这一变化直接改变了算法设计：过去为节省内存而使用的复杂外存策略，在新的容量条件下可能不再必要。

\subsubsection{软件可移植性成为默认期待}

如果相同程序可以在多个处理器代际上继续运行并自动加速，那么软件稳定性和\term{可移植性}{portability}自然具有很高价值。应用代码可以长期演进，而不必每次为了新硬件彻底重写。可扩展且可移植的软件能够跨越多个硬件代际；现代并行软件必须通过更主动的设计才能维持这种寿命。

\subsection{专用计算机在早期为何难以长期保持优势}

专用处理器可以针对某类计算取得初始性能优势，但在通用处理器性能约每 18 个月快速增长的时期，该优势可能在数年内被下一代通用微处理器追平。专用芯片的设计、制造与软件生态成本很高，如果其目标市场不足以支撑持续迭代，就难以与通用处理器的规模经济竞争。

这种现象称为“杀手级微处理器的冲击”\allowbreak (\textit{Attack of the Killer Micros})：许多试图领先于摩尔定律与登纳德缩放组合速度的计算机公司，最终被快速演化的通用微处理器压缩了生存空间。图形处理器之所以比许多其他专用加速器更容易形成长期生态，也与这一背景有关。

\subsection{轻松性能增长的结束}

约在 2005--2006 年，登纳德频率缩放停止延续原有趋势。过去可以忽略的\term{漏电流}{leakage current}变得重要，功耗与散热约束使时钟频率无法继续按历史速度增加。处理器频率大体停留在约 \SIrange{2.5}{3.5}{\giga\hertz} 的范围；近二十年后，频率数量级仍没有发生根本变化。

与此同时，晶体管数量仍在增加，尽管摩尔定律本身也在放缓。于是出现了新的设计矛盾：芯片可以容纳更多晶体管，却不能简单地把它们全部用于提高一个核心的时钟频率。图~\ref{l1:fig:end-easy-gains} 中，晶体管数量继续上升，而频率曲线在 2000 年代中期后明显趋平；逻辑核心数量随后开始增长。

\begin{figure}[H]
  \centering
  \includegraphics[width=0.97\textwidth]{end_easy_performance_gains.png}
  \caption{单纯依靠频率获得性能增长的时代结束}
  \label{l1:fig:end-easy-gains}
\end{figure}

\begin{keybox}
登纳德缩放结束后的关键变化是：硬件仍可能提供更多计算资源，但串行程序不会自动使用这些资源。新增晶体管被组织为更多核心、更多向量执行单元或专用加速器之后，软件必须暴露与之匹配的并行工作，才能把硬件规模转化为应用性能。
\end{keybox}

\subsection{后登纳德时代的三条主要性能路径}

维持性能增长主要有三种方向。

\subsubsection{指令级与向量并行性}

\term{向量执行}{vector execution}让一条指令对多个数据元素执行相同运算。其基本思想是把规则、重复的数据运算打包，使同一控制流驱动多个算术通道。后续图形处理器设计继承并扩大了这种“相同操作作用于不同数据”的思想。

\subsubsection{多核处理器}

当单个核心无法继续大幅提高频率时，可以在同一处理器中放置多个核心。每个核心执行自己的指令流，共享部分缓存、内存与互连资源。多核设计提高了芯片可并行执行的工作量，但单线程程序通常仍只使用一个核心，因此程序需要被改写为多线程或多进程形式。

现代处理器可包含上百个核心。即使硬件核心数从 1 增长到 100，一个只含单一串行指令流的程序也不会自动获得 100 倍加速。核心数量是可用资源，不是无条件性能保证。

\subsubsection{处理器专用化}

通用中央处理器追求“几乎所有代码都能运行，并在多数工作负载上获得尚可到优秀的表现”。专用处理器则主动放弃一部分通用性，把更多芯片面积用于某类操作或数据流。处理器专用化包括以下方向：
\begin{itemize}
  \item 使用更简单的轻量级核心并增加核心数量，例如高性能计算中的蓝色基因系统与至强融核；
  \item 针对向量、流式数据或图形运算构建专用结构；
  \item 在移动芯片中同时配置中央处理器、图形处理器和\term{神经网络处理单元}{Neural Processing Unit, NPU}，分别承担通用控制、图形与特定机器学习运算；
  \item 在超级计算机中由通用中央处理器协调，由图形处理器承担大规模数值计算。
\end{itemize}

\subsection{并行性与异构性的共同转向}

\term{异构计算}{heterogeneous computing}是指同一系统包含多种结构不同、性能特征不同的处理器，并把应用的不同部分映射到最合适的设备。它与并行性相互依赖：专用加速器往往通过大量并行执行单元取得吞吐量优势，而应用必须提供足够规则的并行任务才能利用这些单元。

移动芯片 Apple A14 Bionic 与橡树岭国家实验室 Frontier 百亿亿次级系统都采用加速器架构，说明这种结构从最小规模的移动设备一直延伸到最大规模的超级计算机。

\begin{figure}[H]
  \centering
  \includegraphics[width=0.88\textwidth]{accelerators_across_scales.png}
  \caption{加速器架构贯穿移动设备与超级计算机}
\end{figure}

这一转向改变了软件责任分配。在早期“自动性能增长”环境中，硬件升级可以替软件完成大量加速工作；在后登纳德时代，硬件只提供并行资源与专用能力，程序员必须识别可并行部分、选择可扩展算法、组织数据并消除瓶颈。
